\chapter*{چکیده}
\addcontentsline{toc}{chapter}{چکیده}
\markboth{چکیده}{چکیده}

پیش‌بینی برهم‌کنش‌های زیست‌مولکولی  -  شامل تعاملات دارو-هدف، دارو-دارو، پروتئین-پروتئین و ژن-بیماری  -  یکی از چالش‌های بنیادین در زیست‌شناسی سامانه‌ها و کشف دارو است. مدل \lr{SkipGNN} (هوانگ و همکاران، ۲۰۲۰) با تزریق گراف پرش ۲-گام به شبکه‌های عصبی گراف، عملکرد قابل‌توجهی در بنچمارک‌های استاندارد گزارش کرد. با این حال، بازبینی دقیق تئوریک و تجربی در این پژوهش، چهار محدودیت بنیادین در این مدل شناسایی کرد: بن‌بست هندسی عملگر $A^2$ در گراف‌های دوقطبی، حذف اطلاعات چگالی بر اثر باینری‌سازی، تنگنای ظرفیت دیکودر خطی، و سوگیری ساختاری پروتکل ارزیابی یکنواخت.

برای رفع این محدودیت‌ها، چهار مدل در چارچوب یک پایپ‌لاین یکپارچه طراحی و پیاده‌سازی شدند: \lr{AMS-SkipGNN} (پرش پیوسته تخصیص منابع، گیتینگ تطبیقی، و دیکودر تانسوری ۴‌گانه)، \lr{SkipGATv2} (توجه پویا لبه‌ای تنک)، \lr{ThreeHopSkipGNN} (مسیرهای چندمقیاسی ۱/۲/۳-گام)، و \lr{ContrastiveSkipGNN} (هم‌ترازی خودنظارتی با زیان \lr{InfoNCE}). ارزیابی بر اساس پروتکل دوگانه  -  بانک یکنواخت و بانک سخت آگاه از درجه  -  روی چهار دیتاست استاندارد (\lr{DTI}، \lr{DDI}، \lr{PPI}، \lr{GDI}) با ۳ سید تصادفی انجام شد.

نتایج نشان می‌دهند که در بانک سخت، مدل‌های پیشنهادی بهبود قابل‌توجهی نسبت به \lr{SkipGNN} دارند: در \lr{DDI}، \lr{AMS} به $0.912$ در مقابل $0.724$ می‌رسد؛ در \lr{PPI}، \lr{SkipGATv2} به $0.778$ در مقابل $0.622$؛ و در \lr{DTI} و \lr{GDI}، بهبود در محدوده $+0.10$ تا $+0.12$ واحد درصد است. در دیتاست \lr{DTI}، سه مدل \lr{AMS}، \lr{SkipGATv2} و \lr{3-hop} از نظر آماری هم‌سطح هستند (بازه‌های انحراف معیار هم‌پوشان). در دیتاست \lr{GDI}، هوریستیک تحلیلی تخصیص منابع با $0.842$ بالاتر از تمام مدل‌های عصبی قرار می‌گیرد که به‌عنوان یک بینش مهم درباره محدودیت‌های ذاتی مدل‌های عصبی در گراف‌های بسیار بزرگ و تنک با ویژگی‌های هویتی بحث شده است.

\vspace{1em}
\noindent\textbf{واژگان کلیدی:} پیش‌بینی پیوند، شبکه‌های عصبی گراف، \lr{SkipGNN}، گراف‌های دوقطبی، تخصیص منابع، پروتکل ارزیابی دوگانه
